\documentclass[11pt]{article}
\usepackage[margin=1in]{geometry}
\usepackage{amsmath,amssymb,amsthm,mathtools}
\usepackage{graphicx}
\usepackage{hyperref}
\usepackage{enumitem}
\hypersetup{colorlinks=true,linkcolor=blue,citecolor=blue,urlcolor=blue}
\newtheorem{theorem}{Theorem}
\newtheorem{lemma}[theorem]{Lemma}
\newtheorem{remark}[theorem]{Remark}
\newcommand{\R}{\mathbb R}
\newcommand{\Acal}{\mathcal A}
\newcommand{\PhiA}{\Phi_{\Acal}}
\newcommand{\tr}{\operatorname{tr}}
\newcommand{\op}{\mathrm{op}}
\newcommand{\HS}{\mathrm{HS}}
\newcommand{\Id}{I}

\title{\textbf{Full Solution Packet}\\[3pt]
Alternating operator scaling has the spectral-gap rate}
\author{Run \texttt{fa\_banach\_001}, lane 1}
\date{19 July 2026}

\begin{document}
\maketitle
\begin{center}
\textbf{Status: candidate full solution, likely valid}\\
\textbf{Target: the open question after Corollary 1.6 of arXiv:1904.03213}
\end{center}

\section{Source and exact question}

The source is Kwok--Lau--Ramachandran, \emph{Spectral Analysis of Matrix
Scaling and Operator Scaling}~\cite{KLR}.  For a tuple
$\Acal=(A_1,\ldots,A_k)$, $A_i\in\R^{m\times n}$, put
\[
 s(\Acal)=\sum_i\|A_i\|_F^2,\qquad
 M_{\Acal}=\sum_i A_i\otimes A_i.
\]
It is $\epsilon$-nearly doubly balanced when its two marginals lie between
$1\pm\epsilon$ times their scalar targets, and it has a $\lambda$-spectral
gap when
\[
 \sigma_2(M_{\Acal})\le (1-\lambda)
 \frac{s(\Acal)}{\sqrt{mn}}.
\]
The source proves linear convergence of a gradient flow under
$\lambda^2\ge C\epsilon\log m$ and then asks whether exact alternating left
and right normalization has the same rate.

\begin{center}
\includegraphics[width=.94\linewidth]{figures/open_question_crop.png}
\end{center}

The question is transcribed as follows: ``It is an interesting open question
whether the alternating algorithm also has the same convergence rate as the
gradient flow under the same assumptions.''

\section{Result}

The alternating algorithm replaces $A_i$ successively by
\[
 \left(\sum_j A_jA_j^*\right)^{-1/2}A_i
 \quad\hbox{and}\quad
 A_i\left(\sum_jA_j^*A_j\right)^{-1/2},
\]
with the harmless scalar factors that keep the size normalized.  A full cycle
means one left and one right normalization.

\begin{theorem}[Alternating spectral-gap convergence]\label{thm:main}
There are universal constants $C,c>0$ with the following property.  Let
$m\le n$, $0<\lambda\le1$, and let $\Acal$ be $\epsilon$-nearly doubly
balanced with a $\lambda$-spectral gap.  If
\[
 \lambda^2\ge C\epsilon\log(2m),
\]
then exact alternating normalization is well-defined and, after $t$ full
cycles, has a right-balanced iterate whose other normalized marginal error
$H_t$ satisfies
\[
 \frac{\|H_t\|_F}{\sqrt m}
 \le C\epsilon e^{-c\lambda t}.
\]
Consequently it returns an $\eta$-nearly doubly balanced scaling in
\[
 O\!\left(\frac1\lambda
 \log\frac{m+n}{\eta}\right)
\]
alternating iterations.  Thus the answer to the source question is positive.
\end{theorem}

\section{Proof intuition}

Stop after every right normalization.  Then the right marginal is exactly the
identity and the left marginal is $P=I+H$.  Normalize the completely positive
map to an operator $T$ between Hilbert--Schmidt matrix spaces.  On the scalar
identity directions, $T$ is approximately the identity.  The source's
spectral gap therefore forces the compression of $T$ between the two
traceless subspaces to have norm at most $1-c\lambda$.

Left normalization changes the new right error, to first order, from $H$ to
$-T^*H$; right normalization changes the next left error, to first order, to
$TT^*H$.  Hence one full cycle contracts normalized Frobenius error by
$1-c\lambda$.  The inverse-square-root remainders cost only $O(\epsilon)$.

The remaining issue is global rather than local: could the marginal operator
norm or the spectral gap deteriorate before contraction becomes effective?
For the first $O(\lambda^{-1}\log m)$ cycles, the operator error grows by at
most $1+O(\epsilon)$ per cycle.  The assumed inequality
$\epsilon\log m\ll\lambda^2$ makes the accumulated scaling $O(\lambda)$, so
at least half the gap survives.  By then Frobenius contraction controls the
operator norm directly, and the argument continues to arbitrary accuracy.

\section{One-cycle estimates}

All matrix spaces below carry the Hilbert--Schmidt inner product.  Write
\[
 e_m=I_m/\sqrt m,\qquad e_n=I_n/\sqrt n,
\]
and let $\Pi_m,\Pi_n$ denote orthogonal projection onto the corresponding
traceless subspaces.

Globally rescale so that $s(\Acal)=n$, and suppose first that the right
marginal is exactly balanced:
\[
 \PhiA^*(I_m)=I_n,\qquad
 P:=\frac mn\PhiA(I_n)=I_m+H.
\]
Then $H$ is traceless.  Define
\[
 T=\sqrt{\frac mn}\,\PhiA.
\]
Thus
\[
 T^*e_m=e_n,\qquad Te_n=e_m+\frac{H}{\sqrt m}.
\]

\begin{lemma}[Traceless spectral contraction]\label{lem:trace}
Assume $\|H\|_{\op}\le\delta$, $\sigma_2(T)\le1-\gamma$, and
$\delta\le c_0\gamma$ for a sufficiently small universal $c_0$.  Then
\[
 \|\Pi_mT\Pi_n\|_{\HS\to\HS},\quad
 \|\Pi_nT^*\Pi_m\|_{\HS\to\HS}
 \le 1-c_1\gamma
\]
for a universal $c_1>0$.
\end{lemma}

\begin{proof}
The largest-singular-value estimate from the source~\cite[Lemma 3.6]{KLR}
gives $\sigma_1(T)\le1+\delta$.  Moreover
$\langle e_m,Te_n\rangle=1$.  Apply the elementary spectral lemma
\cite[Lemma 3.7]{KLR}: if unit vectors $p,q$ satisfy
$p^*Tq=1$, $\sigma_1(T)^2\le1+d_1$, and
$\sigma_2(T)^2\le1-d_2$, then
\[
 |x^*Ty|\le1+d_1-d_2
 \quad(x\perp p,\ y\perp q).
\]
Here $d_1\le3\delta$ and $d_2=2\gamma-\gamma^2$.  Reducing $c_0$ if
necessary makes the last bound at most $1-c_1\gamma$.  Taking the supremum
over traceless unit vectors proves both displayed estimates.
\end{proof}

\begin{lemma}[One full normalization cycle]\label{lem:cycle}
In the setting of Lemma~\ref{lem:trace}, put
\[
 \rho=\frac{\|H\|_F}{\sqrt m}.
\]
After exact left normalization followed by exact right normalization, let
$I_m+H_+$ be the new normalized left marginal, and put
$\delta_+=\|H_+\|_{\op}$ and
$\rho_+=\|H_+\|_F/\sqrt m$.  If $\delta\le c_2\gamma$, then
\[
 \rho_+\le(1-c_3\gamma)\rho,
 \qquad
 \delta_+\le(1+C_2\delta)\delta,
\]
for universal constants $c_2,c_3,C_2>0$.
\end{lemma}

\begin{proof}
Left normalization uses $L=P^{-1/2}$.  Its right marginal is $I_n+K$, where
\[
 K=\PhiA^*(P^{-1})-I_n.
\]
It is traceless: using $\PhiA(I_n)=(n/m)P$,
\[
 \tr\PhiA^*(P^{-1})
 =\tr(P^{-1}\PhiA(I_n))=n.
\]
The expansion
\[
 P^{-1}-I_m=-H+H^2P^{-1}
\]
and Lemma~\ref{lem:trace} give
\begin{equation}\label{eq:K}
 \frac{\|K\|_F}{\sqrt n}
 \le (1-c_1\gamma+C_3\delta)\rho.
\end{equation}
Indeed, the traceless part of $T^*H$ contracts by Lemma~\ref{lem:trace}; its
scalar part has norm $\rho\|H\|_F$, because
$\langle e_n,T^*H\rangle=\langle Te_n,H\rangle=\|H\|_F^2/\sqrt m$.
This and the term $H^2P^{-1}$ are both absorbed by $C_3\delta\rho$.

Since $\PhiA^*$ is positive and unital, it is an operator-norm contraction.
Consequently
\begin{equation}\label{eq:Kop}
 \|K\|_{\op}\le\|P^{-1}-I_m\|_{\op}
 \le\frac{\delta}{1-\delta}.
\end{equation}

Now right normalize by $R=(I_n+K)^{-1/2}$.  The next left error is exactly
\[
 H_+=P^{-1/2}\sqrt{\frac mn}\,
 T\big((I_n+K)^{-1}-I_n\big)P^{-1/2}.
\]
Use $(I+K)^{-1}-I=-K+K^2(I+K)^{-1}$.  Because $K$ is traceless,
Lemma~\ref{lem:trace} contracts the linear term.  Equations
\eqref{eq:K}--\eqref{eq:Kop}, together with
$\|P^{-1/2}\|_{\op}^2\le(1-\delta)^{-1}$, therefore yield
\[
 \rho_+\le
 (1+C_4\delta)(1-c_1\gamma+C_4\delta)^2\rho
 \le(1-c_3\gamma)\rho.
\]

For the operator norm, positivity gives
\[
 \left\|\sqrt{\frac mn}T(X)\right\|_{\op}
 \le\|P\|_{\op}\|X\|_{\op}.
\]
Using this with the exact formula for $H_+$ and~\eqref{eq:Kop} gives
\[
 \delta_+
 \le \frac{1+\delta}{1-\delta}
       \frac{\delta}{1-2\delta}
 \le(1+C_2\delta)\delta,
\]
as claimed.
\end{proof}

\section{Gap persistence and proof of Theorem~\ref{thm:main}}

First rescale to size $n$ and perform one right normalization.  Standard
inverse-square-root estimates applied to the initial two marginals give a
right-balanced tuple with
\begin{equation}\label{eq:init}
 \delta_0,\rho_0\le C_5\epsilon,\qquad
 \sigma_2(T_0)\le1-\tfrac34\lambda.
\end{equation}
The last inequality follows because right normalization multiplies
$M_{\Acal}$ on the right by $R\otimes R$ with
$\|R-I\|_{\op}=O(\epsilon)$.  These estimates also show directly that every
normalization used below is invertible.

Let $\delta_j,\rho_j$ be the errors after the $j$th right normalization.
As long as $\delta_j\le2C_5\epsilon$ and the current gap is at least
$\lambda/2$, Lemma~\ref{lem:cycle} gives
\begin{equation}\label{eq:decay}
 \rho_j\le C_5\epsilon e^{-c_3\lambda j},\qquad
 \delta_{j+1}\le(1+C_2\delta_j)\delta_j.
\end{equation}

We record how the gap changes.  A full cycle replaces the natural matrix by
\[
 M\longmapsto(L\otimes L)M(R\otimes R),
\]
where $\|L\|_{\op}^2,\|R\|_{\op}^2\le1+C_6\delta_j$.  Since the size returns
to $n$ after every full cycle,
\begin{equation}\label{eq:gapchange}
 \sigma_2(T_j)\le \sigma_2(T_0)
 \exp\!\left(C_7\sum_{i<j}\delta_i\right).
\end{equation}

Choose
$j_0=C_8\lambda^{-1}\log(2m)$, with $C_8$ large enough that
\eqref{eq:decay} implies $\rho_{j_0}\le C_5\epsilon/\sqrt m$.
On $0\le j\le j_0$, the second inequality in~\eqref{eq:decay} and a bootstrap
give
\[
 \delta_j\le2C_5\epsilon,\qquad
 \sum_{j<j_0}\delta_j
 \le C_9\frac{\epsilon\log(2m)}{\lambda}.
\]
The hypothesis $\lambda^2\ge C\epsilon\log(2m)$, with $C$ sufficiently
large, makes the last quantity at most $c\lambda$.  Equation
\eqref{eq:gapchange} then preserves a gap of at least $\lambda/2$, closing the
bootstrap.

After $j_0$, operator norm is controlled directly by Frobenius norm:
\[
 \delta_j\le\|H_j\|_F=\sqrt m\,\rho_j.
\]
The remaining sum is geometric,
\[
 \sum_{j\ge j_0}\delta_j
 \le \frac{\sqrt m\,\rho_{j_0}}{c_3\lambda}
 \le C_{10}\frac{\epsilon}{\lambda}le c\lambda.
\]
Thus the gap never falls below $\lambda/2$, and Lemma~\ref{lem:cycle} proves
\eqref{eq:decay} for all $j$.

At a right-balanced iterate, the other marginal has operator error at most
$\sqrt m\rho_j$.  Therefore $\sqrt m\rho_j\le\eta$ after
\[
 j=O\!\left(\frac1\lambda
 \log\frac{C\epsilon\sqrt m}{\eta}\right)
 =O\!\left(\frac1\lambda\log\frac{m+n}{\eta}\right)
\]
cycles.  One cycle has two alternating normalizations, so the same asymptotic
bound holds for alternating iterations.  This proves Theorem~\ref{thm:main}.

\section{Verification report}

\begin{itemize}[leftmargin=*]
\item \textbf{Exact target.}  The conclusion is the rate asked for after the
source's Corollary 1.6, under the same near-balance and spectral-gap scale.
\item \textbf{Dimension factors.}  The normalizations
$T=\sqrt{m/n}\,\PhiA$, $P=(m/n)\PhiA(I_n)$, and
$e_d=I_d/\sqrt d$ make the scalar singular directions unit length; all
$\sqrt m,\sqrt n$ factors in the two half-steps cancel.
\item \textbf{Nonlinear terms.}  Both inverse expansions are exact identities,
not first-order approximations.  Their remainders are controlled by the
operator error $\delta$.
\item \textbf{Global bootstrap.}  The only logarithmic dimension loss occurs
before Frobenius error controls operator error.  This is why the source's
$\epsilon\log m\ll\lambda^2$ hypothesis is sufficient.
\item \textbf{Numerical sanity check.}  The included script constructs a
perturbed random balanced quantum expander.  In one fixed run the marginal
error decreased from $6.55\cdot10^{-2}$ to $1.25\cdot10^{-15}$ in 23 cycles;
the spectral gap stayed above $0.4898$, and the median cycle ratio was $0.202$.
The computation is not used in the proof.
\end{itemize}

\section{Novelty and human review recommendation}

The run's registry, solution, attempt, and proof-gap indexes were searched for
arXiv:1904.03213, operator scaling, Sinkhorn, alternating normalization,
spectral gap, and quantum expander; no prior packet was found.  A bounded
external search used the exact question sentence and the phrases ``operator
Sinkhorn spectral gap convergence'' and ``quantum expander alternating
normalization.''  It found the source and the 2024 preprint of
Soma--Uschmajew~\cite{SU}, which proves local asymptotic rates near a fixed
point and studies overrelaxation, but does not state the global theorem above
from a nearly balanced spectral-gap input.  Novelty confidence is therefore
moderate, not absolute.

Send for specialist review.  The highest-value checks are: (i) the use of the
source's spectral lemma as a two-sided traceless compression bound, (ii) the
scalar component of $T^*H$ in~\eqref{eq:K}, and (iii) the cumulative singular-
gap estimate~\eqref{eq:gapchange} for noncommuting successive scalings.

\begin{thebibliography}{9}
\bibitem{KLR}
T. C. Kwok, L. C. Lau, and A. Ramachandran,
\emph{Spectral Analysis of Matrix Scaling and Operator Scaling},
SIAM J. Comput. 50 (2021), no. 3, 1039--1102;
arXiv:1904.03213.

\bibitem{SU}
T. Soma and A. Uschmajew,
\emph{Accelerating Operator Sinkhorn Iteration with Overrelaxation},
arXiv:2410.14104; Math. Program. (2026).
\end{thebibliography}

\end{document}
